
\section{Integrazione e comunicazione}%
\label{sec:integration}

\subsection{Interazione tra sottosistemi e contratti}%
\label{subsec:integration-contracts}

I sottosistemi comunicano attraverso tre protocolli: MQTT tra TMS e CUS, seriale tra WCS e CUS e HTTP tra DBS e CUS.

Nel funzionamento automatico, il TMS invia le misure attraverso il broker MQTT.
Il CUS le elabora, aggiorna lo stato del sistema e, quando necessario, trasmette al WCS il nuovo comando di apertura.
Il WCS aggiorna di conseguenza il servomotore e il display LCD.

Il DBS consulta periodicamente lo stato e lo storico tramite HTTP, indipendentemente dall'aggiornamento del WCS.
La Figura~\ref{fig:sequence} rappresenta il flusso delle informazioni tra i sottosistemi.

\begin{figure}[htbp]
\centering
\begin{tikzpicture}[x=1cm,y=1cm]
\foreach \name/\x in {TMS/0,Broker/3,CUS/6,WCS/9,DBS/12}{
    \node[draw,rounded corners,fill=blue!4,minimum width=1.8cm,
        minimum height=0.7cm] at (\x,0) {\name};
    \draw[dashed,gray] (\x,-0.4) -- (\x,-6.5);
}

\draw[flow] (0,-1) -- node[above,font=\footnotesize]{JSON misura} (3,-1);
\draw[flow] (3,-1.7) -- node[above,font=\footnotesize]{MQTT} (6,-1.7);

\node[font=\footnotesize,align=left,anchor=west] at (6.3,-2.2)
    {storico e politica};

\draw[flow] (6,-3.0) -- node[above,font=\footnotesize]{STATE} (9,-3.0);

\draw[flow] (9,-3.4) -- (9.75,-3.4) -- (9.75,-4.0) -- (9,-4.0);
\node[font=\footnotesize,anchor=west] at (9.9,-3.72)
    {servo e LCD};

\draw[flow] (12,-4.8) -- node[above,font=\footnotesize]{GET stato / storico} (6,-4.8);
\draw[flow] (6,-5.7) -- node[above,font=\footnotesize]{JSON} (12,-5.7);
\end{tikzpicture}
\caption{Sequenza logica di una misura che determina un aggiornamento del sistema. La consultazione del DBS è indipendente dall'attuazione sul WCS}%
\label{fig:sequence}
\end{figure}

\paragraph{MQTT.}

Il TMS pubblica le misure su un topic MQTT condiviso con il CUS, utilizzando messaggi JSON.
Una misura valida contiene lo stato \code{OK}, il livello dell'acqua in metri e la percentuale di riempimento.
Una lettura non valida viene invece segnalata attraverso \code{NO\_ECHO}.

Il CUS utilizza la percentuale ricevuta per la politica automatica e assegna il timestamp ai campioni quando vengono ricevuti.
La comunicazione utilizza QoS 0 e non prevede una conferma applicativa dell'elaborazione della misura.

\paragraph{Seriale.}

La comunicazione tra WCS e CUS utilizza un protocollo testuale a 115200 baud.
Il WCS invia richieste di sincronizzazione e comandi dell'operatore, mentre il CUS trasmette gli aggiornamenti dello stato.

I principali messaggi sono riportati nella Tabella~\ref{tab:serial-protocol}.

\begin{table}[htbp]
\centering
\begin{tabularx}{\linewidth}{@{}l X@{}}
\toprule
\textbf{Messaggio} & \textbf{Significato} \\
\midrule

\code{READY}
& Richiesta di sincronizzazione dal WCS \\

\code{CMD,MODE,\ldots}
& Richiesta di cambio modalità dal WCS \\

\code{CMD,VALVE,\ldots}
& Richiesta di apertura dal WCS \\

\code{STATE,\ldots}
& Stato corrente trasmesso dal CUS al WCS \\

\code{LOG,\ldots}
& Messaggio diagnostico opzionale del WCS \\

\bottomrule
\end{tabularx}
\caption{Principali messaggi del protocollo seriale tra WCS e CUS}%
\label{tab:serial-protocol}
\end{table}

I comandi inviati dal WCS vengono valutati dal CUS, che restituisce lo stato corrente del sistema.
Il messaggio \code{STATE} rappresenta quindi lo stato definito dal controller, non una conferma dell'effettiva posizione del servomotore.

\paragraph{HTTP.}

Il DBS utilizza le API HTTP del CUS per consultare lo stato del sistema, recuperare lo storico delle misure e inviare richieste dell'operatore.

Le operazioni disponibili sono riportate nella Tabella~\ref{tab:http-api}.

\begin{table}[htbp]
\centering
\begin{tabularx}{\linewidth}{@{}l X@{}}
\toprule
\textbf{Operazione} & \textbf{Funzione} \\
\midrule

\code{GET} \path{/api/state}
& Restituisce modalità, apertura e livello corrente. \\

\code{GET} \path{/api/measurements}
& Restituisce gli ultimi \(N\) campioni dello storico. \\

\code{POST} \path{/api/mode}
& Richiede il cambio di modalità operativa. \\

\code{POST} \path{/api/valve}
& Richiede una nuova apertura in modalità MANUAL. \\

\bottomrule
\end{tabularx}
\caption{Operazioni HTTP esposte dal CUS e utilizzate dal DBS}%
\label{tab:http-api}
\end{table}

Le richieste di modifica vengono accettate o rifiutate dal CUS in base allo stato corrente.
Il DBS aggiorna le informazioni visualizzate attraverso successive richieste di lettura dello stato.

\FloatBarrier%

\subsection{Scenari operativi e gestione delle interruzioni}%
\label{subsec:integration-scenarios}

\paragraph{Avvio e sincronizzazione.}

All'avvio il CUS inizializza il sistema in modalità \code{AUTOMATIC}, mentre il WCS mantiene la valvola chiusa fino alla ricezione di uno stato valido.
La sincronizzazione avviene attraverso i messaggi \code{READY} e \code{STATE}.

Il DBS acquisisce autonomamente lo stato tramite HTTP e abilita i controlli in base alla modalità operativa e alla disponibilità del CUS.

\paragraph{Controllo locale e remoto.}

Il WCS e il DBS permettono entrambi di richiedere il cambio di modalità e la regolazione dell'apertura.
Le richieste vengono elaborate dal CUS, senza una priorità prestabilita tra le due interfacce.

Il WCS riceve gli aggiornamenti tramite seriale, mentre il DBS consulta periodicamente lo stato tramite HTTP.
Entrambe le interfacce fanno quindi riferimento alle decisioni del controller.

\paragraph{Indisponibilità delle misure.}

I messaggi \code{NO\_ECHO} rendono il livello corrente non disponibile e azzerano la politica automatica, senza modificare immediatamente l'apertura né rinnovare T2.

In assenza di misure valide per almeno T2, il CUS passa a \code{UNCONNECTED} e imposta l'apertura al \(5\%\).
La condizione viene comunicata al WCS e resa disponibile al DBS.

Alla prima misura valida dopo T2, il CUS ripristina \code{AUTOMATIC} e rivaluta il livello a partire dal fallback interno al \(5\%\).

\paragraph{Interruzioni delle comunicazioni.}

Se il DBS non riesce a comunicare con il CUS, mostra \code{NOT AVAILABLE}, senza modificare lo stato degli altri sottosistemi.

Un'interruzione della comunicazione seriale non impedisce al CUS di elaborare le misure del TMS e le richieste HTTP, ma il WCS può mantenere l'ultimo comando ricevuto.

L'indisponibilità del TMS non implica quindi necessariamente quella del CUS o del DBS.
